\paragraph{Explicit Reasoning Models and Large-scale Reinforcement Learning}

Advances in explicit reasoning models, including OpenAI o1 \& o3~\citep{openai2025o3o4mini, openai2024o1} and DeepSeek R1~\citep{deepseekai2025deepseekr1incentivizingreasoningcapability}, demonstrate their strong ability in mathematical reasoning.
Based on Chain-of-Thought ~\citep{DBLP:journals/corr/abs-2201-11903}, just through prompts, models achieved significant improvements on arithmetic, commonsense, and symbolic reasoning tasks with \textit{explicitly generate} intermediate reasoning text.

Several novel paradigms were also created from explicit reasoning models~\citep{wang2023selfconsistencyimproveschainthought, yao2023treethoughtsdeliberateproblem, Besta_2024}.
In this way, explicit reasoning creates a new level of mathematical reasoning tasks.

\citet{ouyang2022traininglanguagemodelsfollow} showed how GPT-3 can be fine-tuned with reinforcement learning via Proximal Policy Optimization (PPO)~\citep{DBLP:journals/corr/SchulmanWDRK17}, a reinforcement learning algorithm that optimizes the agent's policy by maximizing the expected reward while ensuring that the new policy does not deviate too much from the old policy.
\citet{shao2024deepseekmathpushinglimitsmathematical} proposed Group Relative Policy Optimization (GRPO), which is a new reinforcement learning algorithm that overcame the limitations of PPO, including the need for a large number of samples and the difficulty of tuning hyperparameters.
\citet{yu2025dapoopensourcellmreinforcement} demonstrated the effectiveness of Decoupled Clip and Dynamic Sampling Policy Optimization (DAPO) in large-scale reinforcement learning, which performed best in several benchmarks on Qwen2.5-32B~\citet{qwen2.5} solely based on pure reinforcement learning and the base instruction model of Qwen.

\paragraph{Chain of Continuous Thought}

Previous works have already found that LLMs may involve hidden computation steps, which are not explicitly generated in text~\citep{yang2024largelanguagemodelslatently, chen2023theoremqatheoremdrivenquestionanswering}.
Driven by this spirit, \citet{hao2024traininglargelanguagemodels} proposed \textsc{Coconut}, assigning models to ``reason in a latent continuous space'' through specifically designed training objectives, and demonstrated its potential compared to explicit Chain-of-Thought reasoning procedure, especially on tasks requiring backtracking.
\citet{geiping2025scalingtesttimecomputelatent} introduced a new recurrent transformer architecture that lets models perform multiple hidden reasoning iterations per token.
A recurrent block can be looped arbitrarily at inference to refine the hidden state before producing the next token.
\citet{chen2025language} argued that LLMs may involve a born ability to reason internally, suggesting that LLMs can internally learn better reasoning policies in the latent space when given appropriate training signals.
\citet{saunshi2025reasoninglatentthoughtspower} demonstrated that the critical impact of model capabilities lies in its depth instead of parameter size.
In this way, they proposed \textit{Looped Transformer}, which generated latent thoughts and simulated Chain-of-Thought reasoning procedures internally.

Compared to explicit reasoning, latent reasoning paradigms leverage latent spaces, promoting efficiency and accuracy. However, they still face problems with scaling up and stability.
